\section{Algoritmus \(k\)-NN}\label{sec:ai-knn}

Metoda \(k\) nejbližších sousedů (anglicky \emph{\(k\)-nearest neighbors}, \(k\)-NN) vychází z~intuice, že podobné objekty mívají stejnou třídu. Nevytváří předem parametrický vzorec. Při klasifikaci nového objektu přímo vyhledá nejbližší trénovací objekty a~nechá je hlasovat.

\begin{definition}\label{def:ai-eukleidovska-vzdalenost}
    Pro vektory
    \[
        \mathbf{u}=(u_1,\ldots,u_p),
        \qquad
        \mathbf{v}=(v_1,\ldots,v_p)
    \]
    definujeme \emph{eukleidovskou vzdálenost}
    \[
        d_E(\mathbf{u},\mathbf{v})
        =
        \sqrt{
            \sum_{j=1}^{p}(u_j-v_j)^2
        }.
    \]
\end{definition}

Obecně budeme vzdálenost značit \(d(\mathbf{u},\mathbf{v})\). Seřaďme indexy trénovacích objektů tak, aby
\[
    d(\mathbf{x},\mathbf{x}_{\pi(1)})
    \leqslant
    \cdots
    \leqslant
    d(\mathbf{x},\mathbf{x}_{\pi(n)}).
\]
Množina indexů \(k\) nejbližších sousedů je
\[
    N_k(\mathbf{x})
    =
    \set{\pi(1),\ldots,\pi(k)}.
\]

\begin{definition}\label{def:ai-knn}
    Klasifikátor \(k\)-NN je funkce
    \[
        h^{(k)}(\mathbf{x})
        =
        \argmax_{c\in\mathcal{Y}}
        \sizeof{
            \set{
                i\in N_k(\mathbf{x}):y_i=c
            }
        }.
    \]
\end{definition}

\begin{figure}[ht]
    \centering
    \begin{tikzpicture}
    \begin{axis}[
        width=10cm,height=8cm,
        xmin=0,xmax=8,ymin=0,ymax=6,
        axis lines=left,xlabel={$x_1$},ylabel={$x_2$},
        xtick=\empty,ytick=\empty,
        label style={font=\normalsize},
        every node/.style={font=\normalsize},
        clip=false
    ]
        \addplot[only marks,mark=*,mark size=2.5pt,blue!75!black]
            coordinates{(1,1)(1.8,2)(2.7,1.4)(3.4,2.4)(4,1.7)};
        \addplot[only marks,mark=square*,mark size=2.7pt,orange!85!black]
            coordinates{(4.8,3.3)(5.5,2.8)(5.8,4.2)(6.7,3.5)(7.2,4.8)(6.1,1.9)};
        \addplot[only marks,mark=diamond*,mark size=4pt,red!75!black]
            coordinates{(4.6,2.5)};
        \addplot[only marks,mark=o,mark size=6pt,black]
            coordinates{(4,1.7)(3.4,2.4)(4.8,3.3)(5.5,2.8)(6.1,1.9)};
        \draw[dashed,gray] (axis cs:4.6,2.5)--(axis cs:4,1.7);
        \draw[dashed,gray] (axis cs:4.6,2.5)--(axis cs:3.4,2.4);
        \draw[dashed,gray] (axis cs:4.6,2.5)--(axis cs:4.8,3.3);
        \draw[dashed,gray] (axis cs:4.6,2.5)--(axis cs:5.5,2.8);
        \draw[dashed,gray] (axis cs:4.6,2.5)--(axis cs:6.1,1.9);
        \node[red!75!black,anchor=west] at (axis cs:4.75,2.35){\(\mathbf{x}\)};
    \end{axis}
\end{tikzpicture}

    \caption{Klasifikace nového objektu podle pěti nejbližších sousedů.}
    \label{fig:ai-knn}
\end{figure}

Při shodě hlasů potřebujeme předem stanovené pravidlo. U~binární klasifikace se proto často volí liché \(k\). Malé \(k\) vytváří pružnou hranici citlivou na jednotlivé body a~šum. Velké \(k\) rozhodování vyhlazuje, ale může přehlédnout malé místní skupiny. Hodnotu \(k\) proto volíme podle výsledku na validační množině.

Vedle eukleidovské vzdálenosti lze použít například Minkowského vzdálenost
\[
    d_q(\mathbf{u},\mathbf{v})
    =
    \left(
        \sum_{j=1}^{p}
        \abs{u_j-v_j}^q
    \right)^{1/q},
    \qquad q\geqslant1.
\]
Pro \(q=1\) jde o~manhattanskou vzdálenost, pro \(q=2\) o~eukleidovskou. Smysl výsledku vždy závisí na tom, zda zvolená vzdálenost skutečně vystihuje podobnost objektů dané úlohy.

\importantbox{Před použitím \(k\)-NN je obvykle nutné standardizovat nebo normalizovat jednotlivé vstupní znaky. Jinak může znak s~větším číselným měřítkem zcela ovládnout výpočet vzdálenosti.}

\subsection{Implementace v~Pythonu}

Program~\ref{lst:ai-knn} uchovává trénovací objekty a~jejich třídy. Při predikci spočítá všechny vzdálenosti, vybere \(k\) nejmenších a~vrátí nejčastější třídu. V~případě shody rozhodne menší hodnota třídy, aby byl výsledek jednoznačný.

\lstinputlisting[
    caption={Klasifikátor \(k\)-NN s~eukleidovskou vzdáleností.},
    label={lst:ai-knn}
]{07-ai/code/knn.py}

Trénování je téměř okamžité, ale klasifikace nového objektu vyžaduje porovnání s~celou trénovací množinou. Metoda proto přesouvá většinu výpočetní práce do okamžiku použití modelu.
